% !TeX root = ../../main.tex
\subsection{The two leading coefficients of a single character}
\label{subsec:s05:02}

For one source of degree three, let $T_g^{\rm free}$ and
$T_g^{\rm edge}$ denote the graph sums in
\eqref{eq:s05:source-graphs} with that source assigned, respectively,
to a vertex or to an edge, omitting the Bernoulli prefactor. Thus
\begin{equation}\label{eq:s05:single-source}
 \int_{\MM_g}e^{12\alpha\lambda_1}\mathrm{ch}_3
       =-\frac1{720}(T_g^{\rm free}+T_g^{\rm edge}).
\end{equation}
A free marking gives the leading term. At the next order, terms
from free markings on the $q=1$ cores cancel with the denominator.
A specified self-loop on the $q=0$ core remains. We calculate both
parts before conversion to the normalized Hodge bracket.

\begin{lmm}[One source]\label{lmm:s05:single-source}
As $g\to\infty$,
\begin{align}
 \frac{T_g^{\rm free}}{8\pi^2gZ_g(1)}
 &=f^{(4)}(\alpha)-\frac{5f^{(4)}(\alpha)}{4g}\notag\\
 &\quad-\frac{9f^{(4)}(\alpha)
          +\tfrac{15}{2}\alpha f^{(5)}(\alpha)
          +\alpha^2 f^{(6)}(\alpha)}{\pi^2g}
          +O(g^{-2}),
 \label{eq:s05:free-ratio}\\
 \frac{T_g^{\rm edge}}{Z_g(1)}
 &=-\frac12 f'(\alpha)^2+O(g^{-1}).
 \label{eq:s05:edge-ratio}
\end{align}
\end{lmm}
\begin{prf}
\emph{Truncation before expansion.}
Use \eqref{eq:s05:source-tail} with $r=1$ and $M=2$.
All graphs with no large vertex or with $q\geq2$ contribute
$O(F_g g^{-1})$ to either source sum. On the remaining cores,
\eqref{eq:s05:decorated-core-bound} discards every legal
$p>\log g$ before a large-genus expansion is used. For
$p\leq\log g$ the large genus is $g-O(\log g)$ and its valence is
$O(\log g)$. The HMT thresholds and summable remainder argument in
Proposition~\ref{prp:s03:core-expansion} therefore apply, also after
the fixed derivatives just used in the majorant proof.
Lemma~\ref{lmm:s03:degree-moments} sums their polynomial degree
errors. In this way a relative error $O(g^{-2})$ in a free-marking
term of size $F_g g$ is $O(F_g g^{-1})$ after summation. Only leading
terms are needed for the specified-edge part of size $F_g$.

\emph{The common $q=1$ terms.}
A free marking on a small vertex of a $q=1$ core has $s=0$ in
\eqref{eq:s05:decorated-core-bound}, so it is already
$O(F_g g^{-1})$. A free marking on the large vertex instead
multiplies the leading contribution of that same core by
$8\pi^2g f^{(4)}(\alpha)$. Indeed its leading volume ratio is
$8\pi^2g$, and its differentiated leg factor is
$f^{(4)}(\alpha)$. The core, its small-vertex integrals, its
loop weights and its automorphism factor are unchanged.
This is an additive identity of leading terms for each core and
loop count, without division by a leg factor. Consequently all
$q=1$ contributions to $T_g^{\rm free}/(8\pi^2g)$ equal
$f^{(4)}(\alpha)$ times their denominator contributions at order
$g^{-1}$ and cancel on dividing by $Z_g(1)$. The unmodified
first corrections on the $q=0$ graph cancel in the same way.
We are left with the increment caused by the new marking on that
graph. A specified edge on a $q=1$ core is
$O(F_g g^{-1})$ and needs no coefficient calculation.

\emph{The free-marking increment.}
Write $G=g-p$ and $n=2p$ on the one-vertex graph, and let
$\mathbf d$ be its old degree tuple, with total $D$ and zero count
$m_0$. The first volume coefficient \eqref{eq:s03:c-one} gives
\begin{equation}\label{eq:s05:one-mark-volume}
 \frac{V_{G,n+1}}{8\pi^2G V_{G,n}}
 =1+\frac{n/2-5/4+(2-n)/\pi^2}{G}+O(G^{-2}).
\end{equation}
For fixed $n$ this follows directly by dividing
\eqref{eq:s03:fixed-n-volume}; on the retained growing valences its
coefficients and polynomial remainders have the uniform meaning
established in Section~\ref{sec:s03}. Including $G/g=1-p/g$, the
extra volume factor at order $g^{-1}$ is
$-5/4+(2-2p)/\pi^2$.

The new marking in \eqref{eq:s05:free-markings} has a factor
$\psi_{n+1}^4e^{-\alpha\psi_{n+1}}$. Its descendant degree
$d\geq4$ is therefore always positive, and $m_0$ does not change.
The first profile \eqref{eq:a:first-profile} gives the exact
polynomial difference
\begin{align}
 &e_{n+1,1}(\mathbf d,d)-e_{n,1}(\mathbf d)\notag\\
 &\quad=-\frac{4Dd+2d^2+2D+(2n-3)d-n+2}{\pi^2}.
 \label{eq:s05:profile-increment}
\end{align}
This also holds for the empty old tuple. By
\eqref{eq:a:profile-genus-conversion}, half this difference is the
increment of the coefficient of $G^{-1}$; its source expansion is
in $(2G-2+n)^{-1}$, rather than $G^{-1}$.

For a complex degree-generating variable $z$, put
$\varphi(z)=z^4f^{(4)}(\alpha z)$. Its coefficient of $z^d$
is the weight of the new degree $d$, since differentiating a term
$t^d$ four times leaves degree-counting factor $z^d$.
All these series, also with any fixed degree moment, converge
absolutely. The joint generator of Lemma~\ref{lmm:s04:loop-moments}
gives $\mathcal L[p]=S(1)$, while
\eqref{eq:s04:representative-moments} gives
$\mathcal L[D]=-S(1)$. Substituting $n=2p$ into
\eqref{eq:s05:profile-increment} yields, for each integer $d\geq4$,
\begin{equation}\label{eq:s05:averaged-increment}
 \frac12\mathcal L\left[
   e_{2p+1,1}(\mathbf d,d)-e_{2p,1}(\mathbf d)\right]
 =\frac{2S(1)-d^2+3d/2-1}{\pi^2}.
\end{equation}
The volume increment contributes
$[-5/4+(2-2S(1))/\pi^2]\varphi(1)$.
Adding \eqref{eq:s05:averaged-increment} with weights
$[z^d]\varphi(z)$ cancels $S(1)$. The first correction in the
free ratio is consequently
\begin{equation}\label{eq:s05:free-moment}
 -\frac54\varphi(1)+\frac1{\pi^2}
 \left[\varphi(1)+\left.
  \left(\frac32z\partial_z-(z\partial_z)^2\right)
                  \varphi(z)\right|_{z=1}\right].
\end{equation}
Here
\begin{align*}
 \varphi(1)&=f^{(4)}(\alpha),\\
 (z\partial_z)\varphi(z)|_{z=1}
       &=4f^{(4)}(\alpha)+\alpha f^{(5)}(\alpha),\\
 (z\partial_z)^2\varphi(z)|_{z=1}
       &=16f^{(4)}(\alpha)+9\alpha f^{(5)}(\alpha)
                                     +\alpha^2f^{(6)}(\alpha).
\end{align*}
Substitution proves \eqref{eq:s05:free-ratio}.

\emph{The specified self-loop.}
By \eqref{eq:s05:tagged-edge}, the edge coefficient is
\[
 \mathcal P_3(x,y)e^{-\alpha(x+y)},\qquad
 \mathcal P_3(x,y)=x^2-xy+y^2.
\]
There is no remaining factor $1/(x+y)$ and no integration in the
specified edge parameter.
For $p\geq1$ loops, specifying one of them changes the symmetry
factor by
\begin{equation}\label{eq:s05:specified-loop-weight}
 \frac{p}{2^pp!}=\frac1{2\,2^{p-1}(p-1)!}.
\end{equation}
The other loops sum to $e^{S(1)}$. The specified loop therefore
contributes the factor
\[
 \frac12\left(2f(\alpha)f''(\alpha)-f'(\alpha)^2\right)
       =-\frac12f'(\alpha)^2,
\]
where $f(\alpha)=0$ was used. Its branch-exchange factor $1/2$
has been counted exactly once in
\eqref{eq:s05:specified-loop-weight}. Dividing by the leading
$e^{S(1)}V_{g,0}$ cancels the other loops and the volume factor.
The already controlled remainder is $O(g^{-1})$, proving
\eqref{eq:s05:edge-ratio}.
\end{prf}

Only finitely many derivatives are now needed. For each integer
$k\geq0$ define
\begin{equation}\label{eq:s05:endpoint-derivatives}
 X_k=(8\pi^2)^k f^{(k)}(\alpha).
\end{equation}

\begin{lmm}[Endpoint derivatives]\label{lmm:s05:endpoint-derivatives}
These values satisfy
\begin{equation}\label{eq:s05:derivative-recurrence}
 X_0=0,\quad X_1=-2,\qquad
 X_{k+2}=-(4k+6)X_{k+1}-4\pi^2X_k\quad(k\geq0).
\end{equation}
The required finite list is
\begin{center}
\begin{tabular}{c|l}
$k$ & $X_k$ \\\hline
$0$ & $0$ \\
$1$ & $-2$ \\
$2$ & $12$ \\
$3$ & $8(\pi^2-15)$ \\
$4$ & $1680-160\pi^2$ \\
$5$ & $-32(\pi^4-105\pi^2+945)$ \\
$6$ & $1344(\pi^4-60\pi^2+495)$
\end{tabular}
\end{center}
\end{lmm}
\begin{prf}
The series \eqref{eq:s03:loop-functions} gives
$8tf''(t)+12f'(t)+f(t)=0$. Differentiating $k$ times gives
$8tf^{(k+2)}+(8k+12)f^{(k+1)}+f^{(k)}=0$.
At $t=\alpha=2\pi^2$ this is precisely the recurrence after the
rescaling \eqref{eq:s05:endpoint-derivatives}.
The sine formula gives $f(\alpha)=0$ and
$f'(\alpha)=-1/(4\pi^2)$, hence the two initial values.
Successive substitution gives the table.
\end{prf}

The next calculation includes the factorial conversion in the free
term. Both contributions then have the Hodge bracket normalization.

\begin{prp}[The third character]\label{prp:s05:third-character}
As $g\to\infty$,
\begin{equation}\label{eq:s05:third-character}
 \langle\mathrm{ch}_3\rangle_g
 =\frac{2\pi^2-21}{9g^2}
       +\frac{2(\pi^2-10)}{3g^3}+O(g^{-4}).
\end{equation}
\end{prp}
\begin{prf}
The factor in \eqref{eq:s05:character-normalization} is
\begin{equation}\label{eq:s05:factorial-three}
 (12\alpha)^3\frac{(3g-6)!}{(3g-3)!}
 =\frac{(8\pi^2)^3}{g^3}
             \left(1+\frac4g+O(g^{-2})\right).
\end{equation}
Combining \eqref{eq:s05:single-source} and
\eqref{eq:s05:free-ratio}, the leading coefficient of $g^{-2}$
is $-X_4/720=(2\pi^2-21)/9$.
The free-marking coefficient at order $g^{-3}$, including the
$4/g$ in \eqref{eq:s05:factorial-three}, is
\begin{equation}\label{eq:s05:free-coefficient}
 -\frac1{720}\left[
 \frac{11}{4}X_4-
 \frac{9X_4+15X_5/8+X_6/16}{\pi^2}\right]
 =\frac{29\pi^2}{45}-\frac{20}{3}.
\end{equation}
For the specified loop, \eqref{eq:s05:edge-ratio} gives
\begin{equation}\label{eq:s05:loop-coefficient}
 -\frac{(8\pi^2)^3}{720}
           \left(-\frac12f'(\alpha)^2\right)=\frac{\pi^2}{45}.
\end{equation}
Adding the last two values gives $2(\pi^2-10)/3$.
The $O(g^{-1})$ error before degree conversion, including all graph,
loop and degree tails in Lemma~\ref{lmm:s05:single-source}, becomes
$O(g^{-4})$. The unused second relative term of
\eqref{eq:s05:factorial-three} has the same order after multiplication
by the leading free term. This proves the proposition.
\end{prf}
